\section{Track 综述}

\writingguide{建议写 1200--1500 字。本节必须体现“综述”，不能只罗列论文标题。应按照问题、方法或时间演进建立分类，并引用可靠来源。}

\subsection{研究范围与核心问题}

“\TrackNameEN”关注如何在开放 Web 环境中发现、检索、重组和利用信息。其研究对象既包括传统搜索引擎中的索引、召回、排序和用户交互，也包括面向大语言模型的知识检索、上下文构建、证据归因与生成结果评价。

\placeholder{结合会议 Track 说明补充该方向的正式研究范围，并标注信息来源。}

\subsection{从信息检索到检索增强生成}

传统信息检索通常将系统划分为文档处理、索引构建、候选召回与相关性排序等阶段。神经网络和预训练语言模型推动了稠密检索、跨编码器重排序和端到端问答的发展。RAG 进一步将检索结果作为生成模型的外部上下文，使参数化知识与非参数化知识能够协同工作\parencite{gao2023rag}。

\placeholder{以时间线或技术路线梳理关键阶段，并分析各阶段解决了什么问题、又产生了什么新问题。}

\subsection{典型系统架构}

一个典型的检索增强系统可以抽象为数据准备、查询理解、候选检索、重排序、上下文组织、答案生成和结果评价七个环节。各环节并非彼此独立：分块粒度会影响召回质量，上下文选择会影响生成模型的事实性，而评价指标又决定系统优化的方向。

\placeholder{可在此加入自绘架构图，并分别解释各模块。不要直接复制他人图片；使用他人图片时必须取得许可并注明来源。}
